% WACV 2027 — Apple-to-apple reproduction of PEaRL (arXiv:2510.03455) with
% TabPFN v3 as the prediction head.
%
% \TBD{key} cells are filled by scripts/reproduce_paper.py from JSON outputs.
% Key format:
%   \TBD{table1_gene_<Cohort>_<METRIC>}     ← Table 1 (gene PCC/MSE/MAE)
%   \TBD{table2_pathway_<Cohort>_<METRIC>}  ← Table 2 (pathway PCC/MSE/MAE)
%   \TBD{table3_cindex}                     ← Table 3 (survival C-index)
%
% Run: python scripts/reproduce_paper.py --output-dir wacv_results/paper_run
% Filled output: wacv_results/paper_run/paper.tex
\documentclass[10pt,twocolumn,letterpaper]{article}

\usepackage[utf8]{inputenc}
\usepackage{graphicx}
\usepackage{amsmath}
\usepackage{amssymb}
\usepackage{booktabs}
\usepackage{multirow}
\usepackage{hyperref}
\usepackage{cite}

\title{Apple-to-Apple Reproduction of PEaRL with TabPFN v3 for Gene and Pathway Expression Prediction from Histology}

\author{Ushashi Bhattacharjee \and Alloy Das \and Saria Hannan \and Tirtho Roy \and Soumik Sarkar\\
Iowa State University\\
\texttt{\{ushashi,alloy,saria,tirtho,soumiks\}@iastate.edu}}

\date{}

\begin{document}
\maketitle

\begin{abstract}
We reproduce the PEaRL framework~\cite{majumder2025pearl} for gene and pathway
expression prediction from histology under a strict apple-to-apple protocol on
HEST-1k~\cite{jaume2024hest} and TCGA-BRCA, and replace its MLP prediction
head with TabPFN v3~\cite{priorlabs2024tabpfn}. We characterize accuracy across
three cancer cohorts (Breast, Skin, Lymph), reproduce the published spatial
maps, Leiden clusterings, and survival analysis, and report a head-to-head
comparison of the MLP and TabPFN v3 heads under identical pre-stage-2
encoders. Code and one-command reproduction available at
\url{https://github.com/tirtho149/PeARL-TabFPN}.
\end{abstract}

\section{Introduction}
% (1 page) — frame: PEaRL's contribution + the gap (no in-context head explored),
% TabPFN v3 as a natural plug-in for the MLP, why apple-to-apple matters.
[TODO — write from paper/wacv2027/outline.md \S 1.]

\section{Background}
% (1 page) — minimum needed: pathway-augmented STR prediction, TabPFN v3 in-context
% regression, why a 1-D-target in-context model fits a per-output-dim head bank.
[TODO — see outline.md \S 2.]

\section{Setup}
% (\textonehalf{} page) — frozen task definition. Cohorts, gene targets, pathway
% targets, splits, encoder, seed. The whole point: every comparison row in
% Tables 1--3 uses the same task.
We evaluate on three HEST-1k~\cite{jaume2024hest} cohorts:
Breast (TENX99, 13{,}620 spots, 36 sections, 775 pathways), Skin (TENX158,
8{,}671 spots, 12 sections, 609 pathways), and Lymph node (TENX143, 74{,}220
spots, 24 sections, 1{,}100 pathways). All models train under
\textit{section-stratified GroupKFold 5-fold CV} (matching HEST-Benchmark) with
seed 42, the same data preprocessing pipeline (CP10K $\rightarrow$ log1p $\rightarrow$ 8-neighbor
spatial smoothing $\rightarrow$ per-gene min-max), the same pathway pool (Reactome +
MSigDB Hallmark), and the same UNI v1~\cite{chen2024uni} backbone with the
last four transformer blocks fine-tuned in Stage~1.

\section{Method}
% (\textonehalf{} page) — describe the head swap. MLP $\rightarrow$ TabPFN-v3 (pure mode),
% one regressor per output dim, fit on the trained image embeddings post-Stage-2.
The only thing that differs across head conditions is the post-Stage-2
prediction head. The MLP head is the paper's original two-layer ReLU stack.
The TabPFN v3 head replaces the MLP entirely: one
\texttt{TabPFNRegressor.create\_default\_for\_version(ModelVersion.V3)} per
output dim, fit on the trained image embeddings.

\section{Results}

\subsection{Gene expression prediction (Table~\ref{tab:gene})}

\begin{table}[t]
\centering
\caption{Gene expression prediction (PCC / MSE / MAE; 5-fold mean $\pm$ std). Bold = our run; paper rows reported from arXiv:2510.03455.}
\label{tab:gene}
\begin{tabular}{lccc}
\toprule
Method & Breast & Skin & Lymph \\
\midrule
PEaRL (paper) & $0.5868 \pm 0.0359$ & $0.3756 \pm 0.0148$ & $0.2352 \pm 0.0145$ \\
\textbf{PEaRL+MLP (ours)} & \textbf{\TBD{table1_gene_Breast_PCC}} & \textbf{\TBD{table1_gene_Skin_PCC}} & \textbf{\TBD{table1_gene_Lymph_PCC}} \\
\textbf{PEaRL+TabPFN-v3 (ours)} & \textbf{\TBD{table1_gene_Breast_PCC}} & \textbf{\TBD{table1_gene_Skin_PCC}} & \textbf{\TBD{table1_gene_Lymph_PCC}} \\
\bottomrule
\end{tabular}
\end{table}

\subsection{Pathway expression prediction (Table~\ref{tab:path})}

\begin{table}[t]
\centering
\caption{Pathway expression prediction (PCC / MSE / MAE; 5-fold mean $\pm$ std).}
\label{tab:path}
\begin{tabular}{lccc}
\toprule
Method & Breast & Skin & Lymph \\
\midrule
PEaRL (paper) & $0.5055 \pm 0.0271$ & $0.3523 \pm 0.0336$ & $0.2247 \pm 0.0353$ \\
\textbf{PEaRL+MLP (ours)} & \textbf{\TBD{table2_pathway_Breast_PCC}} & \textbf{\TBD{table2_pathway_Skin_PCC}} & \textbf{\TBD{table2_pathway_Lymph_PCC}} \\
\textbf{PEaRL+TabPFN-v3 (ours)} & \textbf{\TBD{table2_pathway_Breast_PCC}} & \textbf{\TBD{table2_pathway_Skin_PCC}} & \textbf{\TBD{table2_pathway_Lymph_PCC}} \\
\bottomrule
\end{tabular}
\end{table}

\subsection{Survival analysis (Table~\ref{tab:surv})}
% Survival arm — TCGA-BRCA WSI + AB-MIL + Cox partial-likelihood, 5-fold CV.
\begin{table}[t]
\centering
\caption{TCGA-BRCA survival (C-index, 5-fold mean $\pm$ std).}
\label{tab:surv}
\begin{tabular}{lc}
\toprule
Method & C-index \\
\midrule
PEaRL (paper) & $0.659 \pm 0.027$ \\
\textbf{PEaRL+TabPFN-v3 (ours)} & \textbf{\TBD{table3_cindex}} \\
\bottomrule
\end{tabular}
\end{table}

\subsection{Qualitative results}
Spatial heatmaps across the three cohorts for the paper-chosen pathway and gene
per cohort (Figure~\ref{fig:fig3}). Pathway-pathway and gene-gene correlation
matrices (Figures~\ref{fig:fig5},~\ref{fig:fig6}). Leiden clusterings with ARI
labels (Figures~\ref{fig:fig4},~\ref{fig:fig7},~\ref{fig:fig8}).

\begin{figure*}[t]
\centering
\includegraphics[width=0.95\textwidth]{paper_figures/fig3_spatial_predictions.png}
\caption{Spatial predictions of one pathway (top) and one gene (bottom) per cohort.
Breast: Hallmark\_allograft\_rejection / HLA-DMB.
Skin:   Hallmark\_epithelial\_mesenchymal\_transition / QSOX1.
Lymph:  Reactome\_ABC\_family\_proteins\_mediated\_transport / DERL3.}
\label{fig:fig3}
\end{figure*}

\begin{figure}[t] \centering \includegraphics[width=0.95\linewidth]{paper_figures/fig4_leiden_skin.png}
\caption{Leiden clustering on Skin.} \label{fig:fig4} \end{figure}

\begin{figure*}[t] \centering \includegraphics[width=0.9\textwidth]{paper_figures/fig5_pathway_correlation.png}
\caption{Pathway--pathway correlation matrices across the three cohorts.} \label{fig:fig5} \end{figure*}

\begin{figure*}[t] \centering \includegraphics[width=0.9\textwidth]{paper_figures/fig6_gene_correlation.png}
\caption{Gene--gene correlation matrices.} \label{fig:fig6} \end{figure*}

\begin{figure}[t] \centering \includegraphics[width=0.95\linewidth]{paper_figures/fig7_leiden_breast.png}
\caption{Leiden clustering on Breast.} \label{fig:fig7} \end{figure}

\begin{figure}[t] \centering \includegraphics[width=0.95\linewidth]{paper_figures/fig8_leiden_lymph.png}
\caption{Leiden clustering on Lymph.} \label{fig:fig8} \end{figure}

\section{Discussion}
% (\textonehalf{} page) — what changes when we swap MLP $\rightarrow$ TabPFN v3:
% per-cohort delta in PCC, robustness on small cohorts (Skin), failure modes.

\section{Conclusion}
% (\textonefourth{} page) — apple-to-apple holds; TabPFN v3 is a drop-in;
% release the one-command pipeline as a reproducible artifact.

\bibliographystyle{plain}
% Bib will be inlined in the final pass — see paper/wacv2027/refs.bib (TODO).
\begin{thebibliography}{9}

\bibitem{majumder2025pearl}
S.~Majumder, S.~Kapse, M.~Bhattacharya, X.~Xu, A.~Yurovsky, and P.~Prasanna.
\newblock {PEaRL}: Pathway-enhanced representation learning for gene and pathway expression prediction from histology.
\newblock {\em arXiv preprint arXiv:2510.03455}, 2025.

\bibitem{jaume2024hest}
G.~Jaume, P.~Doucet, A.~Song, M.~Y.~Lu, C.~A.~Pérez, et al.
\newblock {HEST-1k}: A dataset for spatial transcriptomics and histology image analysis.
\newblock {\em Advances in Neural Information Processing Systems}, 37:53798--53833, 2025.

\bibitem{chen2024uni}
R.~J.~Chen, T.~Ding, M.~Y.~Lu, et al.
\newblock A general-purpose self-supervised model for computational pathology.
\newblock {\em Nature Medicine}, 30(3):863--874, 2024.

\bibitem{priorlabs2024tabpfn}
N.~Hollmann, S.~Müller, K.~Eggensperger, F.~Hutter.
\newblock {TabPFN}: A transformer that solves small tabular classification problems in a second.
\newblock {\em Nature}, 2024.

\bibitem{ilse2018abmil}
M.~Ilse, J.~Tomczak, M.~Welling.
\newblock Attention-based deep multiple instance learning.
\newblock {\em International Conference on Machine Learning}, 2018.

\end{thebibliography}

\end{document}
